\clearpage\pagestyle{otMain}\setlength{\otBlocEspacement}{2.5mm}\fontsize{11}{15}\selectfont\corrigesChapitre{18}{Algorithmique et programmation par blocs}
\begin{blocLivre}{cor-18-1}\begin{corrige}{1}{Suivre les valeurs}\demarche{Écrire la valeur de x après chaque modification.}
x prend d’abord 4, puis $4+3=7$, puis $7\times2=14$. L’affichage final est donc $\res{14}$. L’ordre des blocs importe.
\label{fig-cor-18-1}\figureGeo[corrige]{script-01}
\end{corrige}\end{blocLivre}
\begin{blocLivre}{cor-18-2}\begin{corrige}{2}{Une entrée utilisateur}\demarche{Traduire la séquence dans son ordre.}
Pour 5, $(5-2)\times4=\res{12}$. Pour une entrée $x$, le programme donne $\res{4(x-2)=4x-8}$. Le dernier bloc utilise la valeur déjà diminuée de 2.
\label{fig-cor-18-2}\figureGeo[corrige]{script-02}
\end{corrige}\end{blocLivre}
\begin{blocLivre}{cor-18-3}\begin{corrige}{3}{Une boucle d’ajouts}\demarche{Compter les tours et situer l’affichage.}
On part de x = 0 et l’on ajoute 5 à chacun des trois tours : 5, 10, 15. Après la boucle, l’affichage est $\res{15}$. Le bloc « dire » est placé une seule fois après la répétition.
\label{fig-cor-18-3}\figureGeo[corrige]{script-03}
\end{corrige}\end{blocLivre}
\begin{blocLivre}{cor-18-4}\begin{corrige}{4}{Une condition au seuil}\demarche{Tester aussi la valeur égale au seuil.}
Pour 9, le test est faux : le message est « petit ». Pour 10, le test est vrai car l’égalité est admise : le message est « grand ». Le cas x = 10 vérifie la frontière de la condition.
\label{fig-cor-18-4}\figureGeo[corrige]{script-04}
\end{corrige}\end{blocLivre}
\begin{blocLivre}{cor-18-5}\begin{corrige}{5}{Un produit répété}\demarche{Distinguer répétition d’un produit et répétition d’une somme.}
Les valeurs après les tours sont 6, 12, 24 et 48. L’affichage est $\res{48}$. On peut écrire $3\times2^4=48$ ; ajouter quatre fois 2 décrirait un autre programme.
\label{fig-cor-18-5}\figureGeo[corrige]{script-05}
\end{corrige}\end{blocLivre}
\begin{blocLivre}{cor-18-6}\begin{corrige}{6}{Une somme construite}\demarche{Suivre les deux variables dans un tableau.}
Au départ s = 0 et n = 1. Les quatre tours ajoutent 1, 2, 3 puis 4 : s vaut successivement 1, 3, 6, 10. Le résultat est $\res{10}$. L’ajout à s précède l’augmentation de n.
\label{fig-cor-18-6}\figureGeo[corrige]{script-06}
\end{corrige}\end{blocLivre}
\begin{blocLivre}{cor-18-7}\begin{corrige}{7}{Tracer un carré}\demarche{Additionner les angles de rotation.}
Il trace quatre segments de 30 pixels, avec un quart de tour antihoraire entre eux. La figure est un carré de côté $\res{30}$ pixels. Il revient en $\res{(0;0)}$ et retrouve sa direction initiale après un tour complet.
\label{fig-cor-18-7}\figureGeo[corrige]{script-07}
\end{corrige}\end{blocLivre}
\begin{blocLivre}{cor-18-8}\begin{corrige}{8}{Tracer un rectangle}\demarche{Lire tout le contenu répété.}
Chaque tour de boucle trace un côté de 50 puis un côté de 20 pixels avec les quarts de tour adaptés. Deux répétitions ferment un rectangle de 50 par 20 pixels. Le périmètre vaut $2(50+20)=\res{140}$ pixels.
\label{fig-cor-18-8}\figureGeo[corrige]{script-08}
\end{corrige}\end{blocLivre}
\begin{blocLivre}{cor-18-9}\begin{corrige}{9}{Deux tests de parité}\demarche{Relier le reste au test de la condition.}
18 a un reste nul dans la division par 2 : le script dit « pair ». 19 a pour reste 1 : il dit « impair ». Les deux branches couvrent ici tous les entiers.
\label{fig-cor-18-9}\figureGeo[corrige]{script-09}
\end{corrige}\end{blocLivre}
\begin{blocLivre}{cor-18-10}\begin{corrige}{10}{Utiliser une touche}\demarche{Distinguer lancement de partie et événement répété.}
Chaque événement « espace pressée » lance un seul ajout. Après cinq appuis, score vaut $\res5$. Mettre score à 0 dans ce même script avant l’ajout donnerait toujours 1 : l’initialisation doit se faire au début de la partie.
\label{fig-cor-18-10}\figureGeo[corrige]{script-10}
\end{corrige}\end{blocLivre}
\begin{blocLivre}{cor-18-11}\begin{corrige}{11}{Un triangle équilatéral}\demarche{Distinguer l’angle intérieur de l’angle de rotation.}
Les trois côtés mesurent 40 et les trois rotations extérieures font $3\times120=360^\circ$. La figure est un triangle équilatéral. L’angle intérieur vaut $180-120=60^\circ$. Tourner seulement de $60^\circ$ ne fermerait pas le triangle en trois côtés.
\label{fig-cor-18-11}\figureGeo[corrige]{script-11}
\end{corrige}\end{blocLivre}
\begin{blocLivre}{cor-18-12}\begin{corrige}{12}{Une interaction simple}\demarche{Repérer le moment où la condition est évaluée.}
Si le contact est présent au moment du test, il dit « Bonjour ! ». Sinon il dit « Je cherche ». Le script ne contient ni répétition ni autre nouveau déclenchement : il effectue un seul test à chaque clic sur le drapeau. Une réaction permanente demanderait de répéter le test.
\label{fig-cor-18-12}\figureGeo[corrige]{script-12}
\end{corrige}\end{blocLivre}
\begin{blocLivre}{cor-18-13}\begin{corrige}{13}{Une dimension saisie}\demarche{Relier variable d’entrée et mesure géométrique.}
La boucle trace un carré de côté c. Pour c = 25, la longueur totale vaut $4\times25=\res{100}$ pixels. En général, elle vaut $\res{4c}$ pixels. Il faut une valeur positive pour la longueur demandée.
\label{fig-cor-18-13}\figureGeo[corrige]{script-13}
\end{corrige}\end{blocLivre}
\begin{blocLivre}{cor-18-14}\begin{corrige}{14}{Chercher une erreur de boucle}\demarche{Conserver la mémoire d’un tour au suivant.}
À chaque tour, s redevient 0 puis vaut 3 : l’affichage est donc $\res3$. Pour accumuler quatre ajouts, déplacer « mettre s à 0 » avant la boucle, puis répéter seulement l’ajout. On obtient $0+4\times3=\res{12}$.
\label{fig-cor-18-14}\figureGeo[corrige]{script-14}
\end{corrige}\end{blocLivre}
\begin{blocLivre}{cor-18-15}\begin{corrige}{15}{Concevoir et vérifier}\demarche{Tester une valeur de chaque côté du seuil et le seuil lui-même.}
Demander x ; mettre y à $3x+5$ ; si y $>0$, dire « positif », sinon dire « nul ou négatif ». Pour 2, y = 11 : positif. Pour $-2$, y = $-1$ : nul ou négatif. Pour $-5/3$, y = 0 : même seconde branche. Le troisième test vérifie l’égalité au seuil.
\label{fig-cor-18-15}\figureGeo[corrige]{script-15}
\end{corrige}\end{blocLivre}
